Recent work trains language agents to predict not only their own actions but also the observations that the environment returns, and reports that this makes later reinforcement learning (RL) more effective because the agent learns the consequences of its actions. We ask whether knowledge learned in this way reaches the agent's decisions. In Know2Act, a pre-registered study of a multi-turn text-to-SQL agent trained with supervised fine-tuning and then GRPO, we measure what each model learns and how it acts. Observation supervision is learned: on unseen databases the agent predicts the format and headers of result tables far better. Yet it changes the agent's actions less than reshuffling the training data does, and it brings no gain after RL. We then vary what the agent predicts and where, up to a contract that states the shape of the correct answer and is checked against every query result. Every prediction is learned, but none changes when the agent submits an answer or keeps querying. The knowledge is usable from outside: selecting among sampled answers with the agent's own contract improves accuracy. RL, which rewards only the final answer, removes even this gain and makes the agent more willing to submit after its own check has failed. Learning to predict the environment is therefore not the same as acting on the prediction. We also find that 44\% of the training tasks in a public SQL RL environment reward an empty answer.
